% =============================================================================
% Section 1 — Introduction
% PINN-ONB01 Phase 1 manuscript (IJHMT submission)
% Target length: ~870 words (6 paragraphs, single flow)
% =============================================================================

\section{Introduction}\label{sec:introduction}

Pool boiling is a passive mechanism capable of dissipating high heat
fluxes. Typical applications span nuclear safety systems, electronics
thermal management, refrigeration evaporators, and process-industry
reboilers.
The onset of nucleate boiling (ONB) marks a regime transition. It is
the wall superheat at which the first sustained vapor embryos detach
from surface cavities, separating inefficient single-phase natural
convection from the high-heat-flux nucleate regime. Accurate knowledge
of $\Delta T_{\mathrm{ONB}} = T_w - T_{\mathrm{sat}}$ is critical. It
sets the lower bound of the design window for evaporators and immersion
cooling. It also defines the safety margin against premature dry-out.
Engineered surfaces now combine roughness control, wettability
tailoring, biphilic patterning, and nano-/micro-structuring;
collectively, these treatments shift $\Delta T_{\mathrm{ONB}}$ by more
than \SI{50}{\percent} relative to a polished
baseline~\citep{betz2013,phan2009,jo2011,bourdon2012,bourdon2015}.
\citet{betz2013} reported an order-of-magnitude reduction in incipience
superheat through biphilic patterning on copper. Bourdon and
co-workers~\citep{bourdon2012,bourdon2015} grafted SAMs on ultra-smooth
bronze ($R_a<\SI{10}{\nano\meter}$) and shifted activation by several
kelvin without modifying topography. Yet design practice still relies
on correlations whose residual scatter remains
\SIrange{20}{30}{\percent}, with consequences that include oversized
exchangers, conservative safety margins, and an inability to exploit
engineered-surface performance.

These persistent design penalties trace back to the structure of the
classical analytical framework itself.
The dominant analytical framework dates to Hsu~\citep{hsu1962}, who
derived an upper- and lower-bound cavity-radius range supporting
nucleation under a linear thermal-boundary-layer assumption. Equivalent
formulations were proposed by Sato and Matsumura~\citep{sato1964}.
Davis and Anderson~\citep{davis1966} derived an explicit closed-form
expression for the minimum superheat, and Bergles and
Rohsenow~\citep{bergles1964} fitted an empirical water-specific
expression with strong predictive power within its calibration band.
Basu et al.~\citep{basu2002} extended the Hsu picture with a
contact-angle modifier $F(\theta)$. These correlations share two
structural limitations. First, each presupposes a single characteristic
cavity radius, precluding multi-modal populations introduced by
engineered patterning. Second, each treats surface character through a
single scalar, most often $R_a$ or $\theta$. Evidence using multiple
descriptors contradicts this
simplification. \citet{jones2009,jabardo2009,betz2013} report
non-additive $\Delta T_{\mathrm{ONB}}$ shifts under joint
roughness--wettability control. As quantified in Section~\ref{sec:results}, these correlations
exhibit RMSEs of \SIrange{7}{16}{K} on a 77-point cross-fluid
benchmark.

These limitations motivate data-driven alternatives that can absorb
multi-descriptor surface information.
Data-driven approaches have been applied to related boiling problems.
Pure data-driven models trained on CHF data have recovered
correlations of comparable accuracy to engineering
fits~\citep{hobold2018,huang2024bubble}, but they remain prone to
unphysical extrapolation and overfitting when surface diversity is
introduced~\citep{li2025cryogenic}. ONB-specific efforts have been
further hindered by data sparsity. Published ONB labels are an order
of magnitude rarer than CHF labels, and they are scattered across
heterogeneous protocols. Black-box training is therefore prone to
overfitting and to violations of thermodynamic trends.
Physics-informed neural networks (PINNs)
\citep{raissi2019,lu2021,cai2021} offer a route around these limits.
They embed governing PDEs and physical constraints as soft losses
during training. This compensates for limited data with mechanistic
regularization and yields models that respect conservation laws under
extrapolation~\citep{zobeiry2021}. Recent work has begun to apply this
machinery to two-phase heat transfer, notably the film-boiling
formulation of \citet{jalili2025pinn} and the physics-assisted CHF
predictors of \citet{li2025cryogenic}. To our knowledge, however, no
PINN study has targeted pool-boiling ONB with explicit
surface-modification conditioning, nor has any integrated curated
multi-source datasets for cross-fluid generalization.

Despite these advances, several concrete gaps remain.
A structured gap analysis identifies four concurrent voids. First, no
published PINN study addresses pool-boiling ONB. Second, none addresses
surface modification, although correlation- and CFD-based baselines
exist for individual surface families. Third, no consolidated
standardized ONB dataset spanning multiple fluids is publicly
available. Fourth, inverse recovery of active cavity radii from
measured ONB data has not been combined with literature-scale
validation against Hsu's canonical bounds~\citep{hsu1962}. Although
\citet{raissi2019} pioneered forward PINN formulations and
\citet{cai2021} extended them to thermal inverse problems, neither
addressed pool-boiling ONB. The present study addresses these four gaps
in a single framework and provides a reproducible baseline for
benchmarking subsequent PINN-for-boiling efforts.

We address these gaps through five contributions.
(i) A curated standardized pool-boiling dataset was assembled from
seven literature sources, comprising $1\,361$ boiling-curve points and
$82$ ONB labels across four fluids (water, R-134a, R-123, FC-77). The
dataset is accompanied by $49$ surface-descriptor cards, of which $37$
are newly digitized; each encodes $R_a$, $\theta$, $r_c$, $N_s$, and a
categorical paper tag.
(ii) A surface-conditioned PINN architecture is proposed. A compact
encoder modulates the temperature-field network through feature-wise
linear modulation (FiLM); the total parameter count is $24\,005$.
One backbone can represent qualitatively distinct surface families.
(iii) A five-term composite loss combines the conduction PDE residual,
the natural-convection boundary condition, the experimental ONB data,
a soft Hsu nucleation criterion, and monotonicity regularizers. The
regularizers enforce
$\partial \Delta T_{\mathrm{ONB}}/\partial R_a \le 0$ and
$\partial \Delta T_{\mathrm{ONB}}/\partial \theta \le 0$.
(iv) An inverse-problem study recovered active cavity radii from ONB
observations using the Hsu relation; the recovered radii were evaluated
against the physical band $r_c \in [1,100]\,\si{\micro\meter}$ and
revealed a Simpson-type reversal in the $R_a$--$r_c$ correlation
across fluids (overall $\rho=-0.27$, water $+0.30$, R-134a $-0.77$).
(v) The framework passed all four Level-1 code-verification tests
(maximum relative error $0.026\%$) and eight of nine Level-3
physical-consistency tests. A deep-ensemble ($K=10$) uncertainty
quantification step is reported. The model attains an ONB RMSE of
\SI{3.42}{K} with $R^2 = +0.44$ on $n=77$ held-out points. This
represents a \SI{53}{\percent} reduction relative to the best classical
correlation. The reduction reaches \SIrange{65}{67}{\percent} on the
refrigerant subset.

The remainder of the paper is organized as follows.
The heat-conduction equation, its natural-convection boundary
condition, the Hsu nucleation criterion, and the non-dimensional scales
are introduced in Section~\ref{sec:formulation}.
Section~\ref{sec:data} then describes the literature sources,
characterization protocol, and preprocessing, while
Section~\ref{sec:architecture} details the PINN architecture, surface
encoder, composite loss, and training. Code verification, forward
accuracy against five baselines, physical-trend consistency, UQ, and
the cavity-radius inverse problem are reported in
Section~\ref{sec:results}. Section~\ref{sec:conclusions} summarizes
findings and outlines the extension to forced-convection boiling.
